\subsection{对与 p-向量相伴的子空间的应用}

设 K 是一个域，E 是 K 上的一个向量空间。回顾一下，对每个 p-向量 \(z \in \bigwedge^{p}(\mathbf{E})\) 相伴着 E 的一个有限维子空间 \(M_{z}\) ，即 E 的使得 \(z \in \bigwedge^{p}(\mathbf{M})\) 的最小向量子空间 M（§ 7，第 2 号，命题 4 的推论）。

命题 14. （i） \(\mathbf{M}_z\) 在 \(\mathbf{E}^*\) 中的正交子空间是 \(x^{*}\in \mathbf{E}^{*}\) （满足 \(x^{*}\lrcorner z = 0\) ）的集合。

\begin{enumerate}
\def\labelenumi{(\roman{enumi})}
\setcounter{enumi}{1}
\tightlist
\item
  子空间 \(\mathbf{M}_z\) （与 \(z\) 相伴）是 \(\bigwedge^{p-1}(\mathbf{E}^*)\) 在映射 \(\lambda_z: u^* \mapsto u^* \lrcorner z\) （从 \(\bigwedge^{p-1}(\mathbf{E}^*)\) 到 \(\mathbf{E}\) ）下的像。
\end{enumerate}

设 \(\mathbf{N}\) 表示 \(\lambda_z\) 的像。对 \(x^* \in \mathbf{E}^*\) 和 \(u^* \in \bigwedge^{p-1}(\mathbf{E}^*)\) ，

\[
\begin{array}{r l} \langle \theta_ {\wedge} (x ^ {*}), u ^ {*} \lrcorner z \rangle = \langle \theta_ {\wedge} (u ^ {*} \wedge x ^ {*}), z \rangle & = (- 1) ^ {p - 1} \langle \theta_ {\wedge} (x ^ {*} \wedge u ^ {*}), z \rangle \\ & = (- 1) ^ {p - 1} \langle \theta_ {\wedge} (u ^ {*}), x ^ {*} \lrcorner z \rangle . \end{array}
\]

因此， \(x^*\) 与 \(\mathbf{N}\) 正交的必要充分条件是 \(x^* \lrcorner z\) 与 \(\theta_\wedge (\bigwedge (\mathrm{E}^*))\) 正交。而后一条件等价于 \(x^* \lrcorner z = 0\) ；因为令 \((e_\lambda)_{\lambda \in L}\) 是 \(\mathbf{E}\) 的一组基；给 \(L\) 一个全序，已知（§ 7，第 8 号，定理 1） \(e_j\) （对 \(J\) 遍历 \(\mathfrak{F}(L)\) （ \(L\) 的有限子集的集合））构成 \(\bigwedge (\mathrm{E})\) 的一组基；于是由第 5 号的公式 (30) 可知，元素 \(\theta_{\wedge}(e_{\mathrm{J}}^{*})\) 是 \(\bigwedge(\mathrm{E})\) 上相对于基 \((e_{\mathrm{J}})\) 的坐标形式（至多差一个符号）；由此得我们的断言。

因此，N 的正交补由 \(x^{*} \in E^{*}\) （满足 \(x^{*} \lrcorner z = 0\) ）组成，从而 (i) 的结论将由 (ii) 得出。

我们首先证明 \(\mathbf{N} \subset \mathbf{M}_z\) 。设 \(\mathbf{M}\) 是 \(\mathbf{E}\) 的一个向量子空间，使得 \(z \in \bigwedge (\mathbf{M})\) ，并设 \(j: \mathbf{M} \to \mathbf{E}\) 是典范单射；令 \(\mu_z\) 表示映射 \(v^* \mapsto v^* \lrcorner z\) （从 \(\bigwedge^{p-1}(\mathbf{M}^*)\) 到 \(\mathbf{M}\) ）；由第 7 号的公式 (60) 可知存在典范分解

\[
\lambda_ {z}: \bigwedge^ {p - 1} (\mathrm{E} ^ {*}) \xrightarrow {\bigwedge^ {p - 1} ({} ^ {t} j)} \bigwedge^ {p - 1} (\mathrm{M} ^ {*}) \xrightarrow {\mu_ {z}} \mathrm{M} \xrightarrow {j} \mathrm{E}
\]

，这证明了 \(\mathbf{N} \subset \mathbf{M}\) ，因此 \(\mathbf{N} \subset \mathbf{M}_z\) （按 \(\mathbf{M}_z\) 的定义）。还需验证 \(\mathbf{N} = \mathbf{M}_z\) 。假设其逆成立：则存在一组基 \((e_i)_{1 \leqslant i \leqslant n}\) （ \(\mathbf{M}_z\) 的基）以及一个元素 \(x^* \in \mathbf{E}^*\) ，使得 \(\langle x^*, e_1 \rangle = 1\) , \(\langle x^*, e_j \rangle = 0\) （对 \(2 \leqslant j \leqslant n\) ），且使得 \(x^*\) 与 \(\mathbf{N}\) 正交，从而 \(x^* \lrcorner z = 0\) 。我们记 \(z = \sum_{\mathbf{H}} a_{\mathbf{H}} e_{\mathbf{H}}\) ，其中求和是对 \([1, n]\) 的含 \(p\) 个元素的子集进行的。由 (68)（第 9 号），

\[
\begin{array}{l l} x ^ {*} \lrcorner e _ {\mathrm{H}} = 0 & \text { if } \quad 1 \notin \mathrm{H} \\ x ^ {*} \lrcorner e _ {\{1 \} \cup \mathrm{H}} = e _ {\mathrm{H}} & \text { if } \quad \mathrm{H} \subset [ 2, n ] \end{array}
\]

这表明关系 \(x^{*} \lrcorner z = 0\) 蕴含 \(a_{\mathbf{H}} = 0\) （对 \(1 \in \mathbf{H}\) ）。但这是不可能的，因为那样 \(z\) 就将属于 \(\bigwedge^{p}(\mathbf{M}')\) ，其中 \(\mathbf{M}'\) 是 \(\mathbf{M}\) 的由 \(e_2, \ldots, e_n\) 生成的子空间。

